\section{Conclusion}
\label{sec:conclusion_v18}

Cartel enforcement often begins before the agency has the
evidence it would need to win a case. Procurement records then
pose an institutional question that is earlier than proof: can
the cheap layer help decide where the expensive layer should be
opened? This paper studies that question on São Paulo's BEC
platform. The answer is yes, but only within a disciplined
legal and empirical scope.

The award-layer statistic ranks firms on persistent zero-win
participation. That ranking is not a membership label. It is a
loser-side exposure measure built from records an agency can
query before recovering bid-level evidence. Validation
therefore targets cartel-adjacent cobidders rather than
direct-defendant status. The ranking survives a conservative
adjudication benchmark, participation-volume discipline,
temporal holdout, leakage audits, and the direct-defendant
scope check that rules out the wrong interpretation of the
screen as a generic cartel-membership detector.

The profile evidence gives the target economic content without
turning it into proof. Cobidders inside the frequent-loser
stratum differ from other frequent losers in deployment
breadth, proximity to direct defendants, product concentration,
and bid-level behavior. These patterns are consistent with the
loser-side screening interpretation and harder to reduce to
mechanical high-volume losing alone. They do not prove cover
bidding, agreement, membership, or liability.

The main contribution is the sequencing result. Bid-distribution
forensics remains necessary, but it need not be the first
stage. Used as a gatekeeper, the award-layer ranking reduces
the bid-microdata pool for the forensic stage by
$\valArchFootprintReduction$ while still recovering
$\valArchTPThouSeqK$ of $\valArchNPos$ adjudicated cobidders.
Joint scoring is the full-observability upper bound; sequential
gatekeeping is the operational architecture for agencies facing
costly bid recovery.

The limits are part of the result. The price evidence is
corroborative, not a damages calculation. The architecture is
portable only where participant identities, winner identities,
repeated comparable procurement, meaningful losing
participation, and recoverable bid files coexist. A static
cutoff can be gamed, so implementation requires refreshed
thresholds, bunching checks, continuous or top-$k$ queues, and
bid-layer follow-up.

Award-layer triage is therefore not a shortcut to liability. It
is a way to order suspicion under incomplete information. The
statistic identifies where legal attention should start, not
where legal responsibility ends.
